\documentclass[border=10pt]{standalone}
\usepackage[no-math]{fontspec}
\usepackage[brazil]{babel}
\usepackage{amsmath,amssymb}
\IfFileExists{C:/Windows/Fonts/Lato_Regular.ttf}{%
  \setmainfont{Lato}[
    Path = C:/Windows/Fonts/,
    UprightFont = Lato_Regular.ttf,
    BoldFont = Lato_Bold.ttf,
    ItalicFont = Lato_Italic.ttf,
    BoldItalicFont = Lato_BoldItalic.ttf
  ]
  \setsansfont{Lato}[
    Path = C:/Windows/Fonts/,
    UprightFont = Lato_Regular.ttf,
    BoldFont = Lato_Bold.ttf,
    ItalicFont = Lato_Italic.ttf,
    BoldItalicFont = Lato_BoldItalic.ttf
  ]
}{%
  \RequirePackage[default]{lato}
}
\renewcommand{\familydefault}{\sfdefault}

\usepackage{xcolor}
\usepackage{tikz}
\usepackage{pgfplots}
\pgfplotsset{compat=1.18}
\usetikzlibrary{arrows.meta,positioning,shapes.geometric,fit,backgrounds,calc,decorations.pathreplacing,patterns}

% Paleta institucional PLANTERR / UEFS (alinhada ao estilo planterr-slides.css / 09a)
\definecolor{etNavy}{HTML}{1F4268}
\definecolor{etNavyDark}{HTML}{152E4A}
\definecolor{etGold}{HTML}{B38A3D}
\definecolor{etGoldDark}{HTML}{8C6927}
\definecolor{etTeal}{HTML}{0E7C7B}
\definecolor{etForest}{HTML}{2D6A4F}
\definecolor{etAlert}{HTML}{9E2A2B}
\definecolor{etInk}{HTML}{26313D}
\definecolor{etMuted}{HTML}{5B6573}
\definecolor{etBorder}{HTML}{D5DEE8}

\definecolor{bgNavy}{HTML}{EDF2F7}
\definecolor{bgGold}{HTML}{F8F3E9}
\definecolor{bgTeal}{HTML}{EBF6F5}
\definecolor{bgForest}{HTML}{EEF5F0}
\definecolor{bgAlert}{HTML}{FDF2F2}
\definecolor{bgCanvas}{HTML}{FFFFFF}

\tikzset{
  >=Stealth,
  etBox/.style={
    draw=#1, fill=white, rounded corners=5pt, line width=1.15pt,
    inner sep=7pt, align=center, font=\footnotesize\sffamily, text=etInk
  },
  etTint/.style={
    draw=#1, fill=#1!8!white, rounded corners=5pt, line width=1.15pt,
    inner sep=7pt, align=center, font=\footnotesize\sffamily, text=etInk
  },
  etBadge/.style={
    fill=#1, text=white, rounded corners=9pt, inner xsep=8pt, inner ysep=3pt,
    font=\scriptsize\bfseries\sffamily
  },
  etTitle/.style={font=\normalsize\bfseries\sffamily, text=etNavy, align=center},
  etSub/.style={font=\footnotesize\sffamily, text=etMuted, align=center},
  etArrow/.style={->, >=Stealth, line width=1.35pt},
  etDArrow/.style={<->, >=Stealth, line width=1.35pt}
}


\begin{document}
\begin{tikzpicture}[font=\sffamily]

\begin{axis}[
  width=14.2cm,
  height=8.4cm,
  title={\large\bfseries\sffamily\color{etNavy} Perfil do VPL $\times$ Taxa de Desconto (Conflito VPL \textit{vs.} TIR)},
  title style={yshift=6pt},
  xlabel={Taxa de Desconto (\% ao ano)},
  ylabel={Valor Presente Líquido (R\$ mil)},
  label style={font=\small\bfseries\sffamily, text=etInk},
  tick label style={font=\footnotesize\sffamily, text=etMuted},
  xmin=0, xmax=30,
  ymin=-120, ymax=460,
  xtick={0,5,10,15,20,25,30},
  ytick={-100,0,100,200,300,400},
  grid=major,
  grid style={color=etBorder!80, line width=0.45pt},
  axis line style={color=etInk!70, line width=0.8pt},
  tick style={color=etInk!60},
  legend style={
    at={(0.97,0.97)}, anchor=north east,
    draw=etBorder, fill=white, fill opacity=0.95, text opacity=1,
    rounded corners=4pt, inner sep=5pt,
    font=\scriptsize\sffamily, cells={anchor=west}
  },
  clip=false
]

% Linha horizontal zero (eixo de decisão VPL = 0)
\draw[etInk!85, line width=1.0pt] (axis cs:0,0) -- (axis cs:30,0);

% Linha vertical da TMA = 10% a.a.
\draw[etForest, dotted, line width=1.6pt] (axis cs:10,-120) -- (axis cs:10,450);

% Curva Projeto A (Maior escala)
\addplot[
  color=etNavy, line width=2.2pt, smooth
] coordinates {
  (0,420) (5,286) (10,174) (16,74) (21.8,0) (30,-94)
};
\addlegendentry{\textbf{Projeto A} — Maior Escala (VPL maior para TMA $< 16\%$)}

% Curva Projeto B (Menor escala, TIR maior)
\addplot[
  color=etTeal, dashed, line width=2.2pt, smooth
] coordinates {
  (0,240) (5,176) (10,122) (16,74) (26.8,0) (30,-16)
};
\addlegendentry{\textbf{Projeto B} — Menor Escala (TIR maior)}

% Pontos marcadores principais
\addplot[only marks, mark=*, mark size=3.2pt, color=etGoldDark, fill=etGold] coordinates {(16,74)};
\addplot[only marks, mark=*, mark size=3.2pt, color=etNavy, fill=etNavy] coordinates {(21.8,0)};
\addplot[only marks, mark=*, mark size=3.2pt, color=etTeal, fill=etTeal] coordinates {(26.8,0)};
\addplot[only marks, mark=*, mark size=2.8pt, color=etForest, fill=etForest] coordinates {(10,174) (10,122)};

% Callout TMA = 10% a.a. (posicionado à esquerda da reta vertical, livre da legenda)
\node[
  draw=etForest, fill=bgForest, rounded corners=4pt, inner sep=5pt,
  font=\scriptsize\bfseries\sffamily, text=etForest, align=center
] (tma_box) at (axis cs:5.1, 345) {TMA $= 10\%$ a.a.\\[-1pt]\normalfont\scriptsize Aceitar se $\text{VPL} > 0$};
\draw[->, >=Stealth, line width=1.1pt, color=etForest] (tma_box.east) -- (axis cs:9.8, 325);

% Callout Taxa de Fisher (~16%)
\node[
  draw=etGold, fill=bgGold, rounded corners=4pt, inner sep=5pt,
  font=\scriptsize\bfseries\sffamily, text=etGoldDark, align=center
] (fisher_box) at (axis cs:19.2, 185) {Taxa de Fisher $\approx 16\%$\\[-1pt]\normalfont\scriptsize $\text{VPL}_A = \text{VPL}_B$};
\draw[->, >=Stealth, line width=1.1pt, color=etGoldDark] (fisher_box.south west) -- (axis cs:16.3, 81);

% Callout TIR A (~21,8%)
\node[
  draw=etNavy, fill=bgNavy, rounded corners=4pt, inner sep=4.5pt,
  font=\scriptsize\bfseries\sffamily, text=etNavy, align=center
] (tira_box) at (axis cs:17.2, -62) {$\text{TIR}_A \approx 21,8\%$};
\draw[->, >=Stealth, line width=1.1pt, color=etNavy] (tira_box.north east) -- (axis cs:21.4, -7);

% Callout TIR B (~26,8%)
\node[
  draw=etTeal, fill=bgTeal, rounded corners=4pt, inner sep=4.5pt,
  font=\scriptsize\bfseries\sffamily, text=etTeal, align=center
] (tirb_box) at (axis cs:26.0, 92) {$\text{TIR}_B \approx 26,8\%$};
\draw[->, >=Stealth, line width=1.1pt, color=etTeal] (tirb_box.south) -- (axis cs:26.7, 9);

\end{axis}
\end{tikzpicture}
\end{document}
